% SPDX-License-Identifier: Apache-2.0
\section{Numbers Nobody Can Predict}
\label{sec:x-trng}

A network stack wants numbers nobody can guess: the first sequence
number of a connection, a nonce, a key. A machine with no clock of the
day, which boots into the same state every time, has nowhere to get
them but physics, and until this part the Vreteno machine had nowhere
at all: Zephyr's build with a network stack took a counter and called
it random (issue 458).

The physics is the jitter of a ring oscillator. A loop of an odd
number of inverters runs at a rate set by its own gates, wandering
with temperature and supply, and where the clock's edge catches it in
its period is not something a program can know. Eight such rings are
\code{RingOsc}, a Verilog module the netlist instantiates and does not
write, since a loop of gates is not a thing this language says; the
attributes on it stop Vivado folding the loops to nothing. In a Rust
run it is a model, a shift register with feedback per ring, and the
example says so: nothing it prints is evidence of randomness, only of
the machinery. Only a board is evidence of that, and the issue says
how far it has been measured.

\code{Trng} is the other half and is ordinary hardware. Each cycle it
folds the eight samples to one bit by exclusive or, which is less
biased than any one of them, and does three things with the stream.
It watches it: forty identical bits in a row, the repetition count
test of NIST SP 800-90B at a false alarm rate of one in a million,
raise a sticky fault and stop the buffer filling, since a ring that
has stopped is worse than none. It debiases it: von Neumann's
extractor takes the bits in pairs, \code{01} a zero, \code{10} a one,
and the two equal pairs nothing, so whatever a ring's bias is, it
cancels, at the cost of three bits in four. And it buffers it: the
bits that survive fill a word, and whole words wait in a buffer of
four, which a host takes one at a time from \code{data} after
\code{status} says one is ready. A fourth word, \code{raw}, is the
last thirty-two samples before the extractor, and is what a program
on the board reads in a loop to measure the source.

\code{Entropy} is the two joined, the samples one way and the run bit
the other, so that a design holds one unit and its netlist holds the
module inside.

\begin{figure}[htbp]
\centering
\input{trng_timing.tex}
\caption{The source starting: \code{en} rises with the run bit, the
samples arrive on \code{raw}, \code{have} and \code{first} hold the
first bit of each pair, \code{nbits} climbs by the pairs that
differ, and \code{runlen} counts the identical samples the health
test watches.}
\label{fig:x-trng}
\end{figure}

\lstinputlisting[caption={\code{ex\_trng.rs}. Run with \code{bazel run //lib/examples:ex\_trng}.},label={lst:x-trng}]{lib/examples/ex_trng.rs}

\lstinputlisting[style=ascii,caption={What \code{ex\_trng} printed when the build ran it: the eight words, the raw samples, and the buffer drained after the stop.},label={lst:x-trng-out}]{out_trng.txt}
